\documentclass{article}
\usepackage[a4paper, margin=1in]{geometry}

\usepackage{amsmath, amssymb}
\usepackage{setspace}
\setstretch{1.2}
\setlength{\parindent}{0.5em}
\setlength{\parskip}{0.5em}
\usepackage{lmodern}
\usepackage{amsmath}
\usepackage{enumitem}
\usepackage{hyperref}


\begin{document}


\title{Evaluating if a strategy pair is SPNE in the EcoPG stochastic game}
\maketitle


\section{Definitions}

To evaluate if a strategy pair is SPNE, we need to check if any agent has an incentive to make a one-step deviation from the strategy and then return 
to the same strategy, given the other agent is following the strategy. If there is no incentive to deviate for any agent, then the strategy pair is SPNE.

By incentive to deviate, we mean that the cumulative reward for the agent from deviating for a single step and then returning to the strategy is greater 
than the cumulative reward from following the strategy without deviation.

To estimate the cumulative rewards in either case (i.e., with deviation and without deviation), 
we write value functions for the agents. 

In the EcoPG, the agents start in the prosperous state, and can transition to the degraded state with a probability dependent on the joint action. 
Once in the degraded state, they can either stay there or recover back to the prosperous state with some probability again dependent on their joint actions.

To derive the value functions for the prosperous state $V_P$ and the degraded state $V_D$, we can use the Bellman equations.

Let the immediate payoff for the agent in the prosperous state be $u$, and the penalty associated with transitioning (or staying) in the degraded state be $m$.
 $\gamma$ is the discount factor, $p_c$ is the probability of collapse from the prosperous state to the degraded state, and $p_r$ is the probability of recovery from the degraded state to the prosperous state.

Note that $u$, $p_c$, and $p_r$ are all functions of the joint action of the agents, and hence should be written as $u(a_1, a_2)$, $p_c(a_1, a_2)$, and $p_r(a_1, a_2)$
 respectively.However, for notational simplicity, we will write them without the arguments in the derivations below.

% V_P definition
The value function in the prosperous state is
\begin{equation}
\label{eq:VP_def}
\begin{aligned}
V_{P} &= \left(1 - p_{c}\right) \left(V_{P} \gamma + u\right) +  p_{c} \left(V_{D} \gamma + m\right)
\end{aligned}
\end{equation}

Basically, 1 - $p_c$ is the probability of staying in the prosperous state, in which case the immediate payoff is $u$ and the discounted future payoff is $\gamma V_P$.
$p_c$ is the probability of transitioning to the degraded state, in which case the immediate payoff is $m$ and the discounted future payoff is $\gamma V_D$.

% V_D definition and rearrangement
Similarly, the value function in the degraded state is
\begin{equation}
\label{eq:VD_def}
\begin{aligned}
V_{D} &=  m  + \gamma \left(V_{D} \left(1 - p_{r}\right) + V_{P} p_{r}\right) \\
% V_{D} &= \frac{V_{P} \gamma p_{r} + m}{\gamma p_{r} - \gamma + 1}
\end{aligned}
\end{equation}

The immediate payoff will always be $m$ in the degraded state, and the future payoff is a weighted average of the payoffs from staying in the degraded state (with probability $1 - p_r$) and recovering to the prosperous state (with probability $p_r$).

Rearraging the above equation, we get an expression for $V_D$ 
\begin{equation}
\label{eq:VD_def_rearranged}
\begin{aligned}
V_{D} &= \frac{V_{P} \gamma p_{r} + m}{\gamma p_{r} - \gamma + 1}
\end{aligned}
\end{equation}


Now, \ref{eq:VD_def_rearranged} can be substituted into \ref{eq:VP_def} to get an expression for $V_P$ in terms of the parameters of the model.
Doing so and simplifying yields

\begin{equation}
\label{eq:VP_final}
\begin{aligned}
V_{P} &= \left(1 - p_{c}\right) \left(V_{P} \gamma + u\right) + p_{c} \left(\frac{\gamma \left(V_{P} \gamma p_{r} + m\right)}{\gamma p_{r} - \gamma + 1} + m\right) \\
V_{P} &= \frac{\gamma m p_{c} p_{r} - \gamma p_{c} p_{r} u + \gamma p_{c} u + \gamma p_{r} u - \gamma u + m p_{c} - p_{c} u + u}{- \gamma^{2} p_{c} - \gamma^{2} p_{r} + \gamma^{2} + \gamma p_{c} + \gamma p_{r} - 2 \gamma + 1} \\
V_{P} &= \frac{m p_{c} \left(1 + \gamma p_{r}\right) + u \left(1 - p_{c}\right) \left(1 - \gamma \left(1 - p_{r}\right)\right)}{\left(1 - \gamma\right) \left(1 + \gamma p_{r} - \gamma \left(1 - p_{c}\right) \right)}
\end{aligned}
\end{equation}

We will use the last equation in all subsequent steps.

\begin{equation}
    V_{P} = \frac{m p_{c} \left(1 + \gamma p_{r}\right) + u \left(1 - p_{c}\right) \left(1 - \gamma \left(1 - p_{r}\right)\right)}{\left(1 - \gamma\right) \left(1 - \gamma \left(1 - p_{c}\right) + \gamma p_{r} \right)}
\end{equation}

\section{Comparison of Cumulative Rewards with and without Deviation}

Now, we compare the cumulative rewards for the agent from deviating for a single step and then returning to the strategy, 
with the cumulative reward from following the strategy without deviation.


Without deviation, the cumulative reward is simply $V_P$ as derived above.

With deviation, we need to consider the consequences of a single step deviation. Let us say the agent deviates at time step $t=0$. 
The immediate payoff from deviation is $u^0$, and the probability of collapse due to deviation is $p^0$.

So, the cumulative reward from deviation at $t=0$ can be written as:

\begin{equation}\label{eq:VP_deviation}
    V_P^{(0)} = (1-p^0)(u^0 + \gamma V_P^{(1)}) + p^0(m + \gamma V_D^{(1)})
\end{equation}

In \ref{eq:VP_deviation}, $V_P^{(1)}$ is the cumulative reward from time step $t=1$ onwards, given that the agent deviated at $t=0$ and then returned to the strategy.
Likewise, $V_D^{(1)}$ is the cumulative reward from time step $t=1$ onwards in the degraded state, given that the agent deviated at $t=0$ and then returned to the strategy.
The above equation is the general form of the SPNE equation
Thus, if a strategy is SPNE, then the following condition must hold:

\begin{equation}\label{eq:SPNE_inequality}
    V_P^{(0)} \leq V_P 
\end{equation}

\begin{equation}
\boxed{
(1-p^{(0)})(u^0 + \gamma V_P^{(1)}) + p^{(0)}(m + \gamma V_D^{(1)})
\leq
\frac{m p_{c} \left(1 + \gamma p_{r}\right) + u \left(1 - p_{c}\right) \left(1 - \gamma \left(1 - p_{r}\right)\right)}
{\left(1 - \gamma\right) \left(1 + \gamma p_{r} - \gamma \left(1 - p_{c}\right) \right)}
}
\end{equation}

\subsection{SPNE inequality in different info conditions}

We look at how how $V_P^{(0)}$ differs across the different information conditions.

\subsubsection{Only Ecological Information: }

Under the only ecological information condition, the agents can distinguish betweeen the prosperous and dedgraded states, but they play memory-0 strategies.
So, a one step deviation will affect only the payoff and the probability of collapse at $t=0$, but this effect does not carry on the future steps.

So, we can write the cumulative reward from deviation as:
\begin{equation}
    V_P^{(0)} = (1-p^{(0)})(u^0 + \gamma V_P) + p^{(0)}(m + \gamma V_D) 
\end{equation}

This is simply \ref{eq:VP_deviation} with $V_P^{1} = V_P$ and $V_D^{1} = V_D$.

\subsubsection{Complete Information}:

If agents play memory-1 strategies in addition to having to knowledge of the ecological state, 
then a one step deviation can have consequences beyond just the time-step in which the deviation occurs.

For example, consider a strategy like WSLS. If both agents are playing WSLS and cooperating from the beginning, then if one agent deviates at $t=0$ by defecting, then at $t=1$, both agents will defect.
However, at $t=2$, both agents will again cooperate. So, the effect of deviation lasts for one more time step after the deviation step. 
So, we need to compute $V_P^{1}$ and $V_D^{1}$ separately, since they will differ from $V_P$ and $V_D$ respectively. But $V_P^{2}$ and $V_D^{2}$ will be the same as $V_P$ and $V_D$ respectively, since after $t=2$, the effect of deviation is gone and both agents are back to cooperating.

\begin{equation}\label{V_P_deviation_with_t_1}
\begin{aligned}
    V_{P}^{(0)} &= \left(1 - p_{c}^{(0)}\right)\left( u + \gamma V_{P}^{(1)} \right) 
    + p_{c}^{(0)}\left( m + \gamma V_{D}^{(1)} \right)
\end{aligned}
\end{equation}

\begin{equation}\label{V_P_from_t_1_onwards}
\begin{aligned}
    V_{P}^{(1)} &= \left(1 - p_{c}^{(1)}\right)\left( u + \gamma V_{P}^{(1)} \right) 
    + p_{c}^{(1)}\left( m + \gamma V_{D}^{(1)} \right)\\
    V_{D}^{(1)} &=  m + \gamma \left( (1 - p_{r}^{(1)}) V_{D}^{(1)} + p_{r}^{(1)} V_{P}^{(1)} \right)
\end{aligned}
\end{equation}



If needed in GT like strategies, the effect of deviation can be permanent, and the agents that initially cooperated keep defecting forever after the deviation. In this case,
we need to compute $V_P^{(2)}$ and $V_D^{(2)}$ separately as well. However, after $t=2$, the effect of deviation is permanent, so $V_P^{(3)} = V_P^{(2)}$ and $V_D^{(3)} = V_D^{(2)}$.


\begin{equation}
\begin{aligned}
    V_{P}^{(0)} &= \left(1 - p_{c}^{(0)}\right)\left( u + \gamma V_{P}^{(1)} \right) 
    + p_{c}^{(0)}\left( m + \gamma V_{D}^{(1)} \right)
\end{aligned}
\end{equation}

\begin{equation}
\begin{aligned}
    V_{P}^{(1)} &= \left(1 - p_{c}^{(1)}\right)\left( u + \gamma V_{P}^{(1)} \right) 
    + p_{c}^{(1)}\left( m + \gamma V_{D}^{(1)} \right)\\
    V_{D}^{(1)} &=  m + \gamma \left( (1 - p_{r}^{(1)}) V_{D}^{(1)} + p_{r}^{(1)} V_{P}^{(1)} \right)
\end{aligned}
\end{equation}

\begin{equation}
\begin{aligned}
    V_{P}^{(2)} &= \left(1 - p_{c}^{(2)}\right)\left( u + \gamma V_{P}^{(2)} \right) 
    + p_{c}^{(2)}\left( m + \gamma V_{D}^{(2)} \right)\\
    V_{D}^{(2)} &=  m + \gamma \left( (1 - p_{r}^{(2)}) V_{D}^{(2)} + p_{r}^{(2)} V_{P}^{(2)} \right)
\end{aligned}
\end{equation}


\subsubsection{No Information}

Under the no information condition, the agents cannot distinguish between the prosperous and degraded states. 
So, from the POV of the agents, there is no $V_P$ and $V_D$, but a single value function $V_{\text{no info}}$ which is the same for both states.

\textbf{WORK IN PROGRESS}

\subsubsection{Only Social Information}

\textbf{WORK IN PROGRESS}





\end{document}


